Write \(M(n,k)=\operatorname{lcm}(n+1,\ldots,n+k)\).
There are infinitely many triples \(n,m,k\), with \(k\geq3\) and
\(m\geq n+k\), for which \(M(n,k)>M(m,k+1)\).
In fact, for every fixed \(C\geq1\) and all sufficiently large \(k\)
we construct disjoint intervals beginning at \(x<y\) such that
\[
 \operatorname{lcm}(x,\ldots,x+k-1)>
 C\operatorname{lcm}(y,\ldots,y+k).
\]
This reconstructs Cambie's theorem, with a simplified choice of \(y\):
\url{https://arxiv.org/abs/2410.09138},
\url{https://www.erdosproblems.com/678}.
The only asymptotic external input is the prime number theorem, used
to find primes in fixed relative intervals.
